% Included inside the existing thebibliography environment.
% Verification limits: Sabbagh and Law claims use publisher abstracts;
% Sarda/Lena uses the publisher's chapter outline; Kimoto uses the abstract.
% Latrouite/Van Valin full primary PDF checked, especially pp. 161--170.
% Wellstood's full text was blocked by publisher site protection; the relative-
% clause/possessor-raising proposal is confirmed by the indexed primary abstract.
% Consultant counts and specific extraction judgments are not asserted here.
\bibitem{sabbagh2009}
Sabbagh, Joseph. 2009.
Existential sentences in Tagalog.
\textit{Natural Language \& Linguistic Theory} 27(4): 675--719.
\href{https://doi.org/10.1007/s11049-009-9083-3}{doi:10.1007/s11049-009-9083-3}.

\bibitem{aldridge2011event}
Aldridge, Edith. 2011.
Event existentials in Tagalog.
In Lauren Eby Clemens, Gregory Scontras, and Maria Polinsky (eds.),
\textit{Proceedings of the Eighteenth Meeting of the Austronesian Formal
Linguistics Association}, 16--30.
\href{https://ojs.lib.uwo.ca/index.php/afla/issue/view/2037}{University of Western Ontario proceedings archive}.
The archive dates the volume 2011; the author's University of Washington
publication record dates the proceedings contribution 2012.

\bibitem{law2010impersonal}
Law, Paul. 2010.
The impersonal construction in Tagalog.
In Raphael Mercado, Eric Potsdam, and Lisa deMena Travis (eds.),
\textit{Austronesian and Theoretical Linguistics}, 297--326.
Amsterdam: John Benjamins.
\href{https://doi.org/10.1075/la.167.17law}{doi:10.1075/la.167.17law}.

\bibitem{latrouitevanvalin2014event}
Latrouite, Anja, and Robert D. Van Valin, Jr. 2014.
Event existentials in Tagalog: A Role and Reference Grammar account.
In I Wayan Arka and Ni Luh Ketut Mas Indrawati (eds.),
\textit{Argument Realisations and Related Constructions in Austronesian
Languages: Papers from 12-ICAL}, vol.~2, 161--174.
Canberra: Asia-Pacific Linguistics.
\href{https://rrg.caset.buffalo.edu/rrg/vanvalin_papers/Lat_Valin.pdf}{Author-hosted chapter}.

\bibitem{wellstood2022}
Wellstood, Zachary. 2022.
A relative clause analysis of event existential constructions in Aklanon.
\textit{Glossa: A Journal of General Linguistics} 7(1): 1--41.
\href{https://doi.org/10.16995/glossa.5866}{doi:10.16995/glossa.5866}.

\bibitem{mcnally2016existential}
McNally, Louise. 2016.
Existential sentences crosslinguistically: Variations in form and meaning.
\textit{Annual Review of Linguistics} 2: 211--231.
\href{https://doi.org/10.1146/annurev-linguistics-011415-040837}{doi:10.1146/annurev-linguistics-011415-040837}.

\bibitem{sardalena2023}
Sarda, Laure, and Ludovica Lena. 2023.
Existential constructions: In search of a definition.
In Laure Sarda and Ludovica Lena (eds.),
\textit{Existential Constructions across Languages: Forms, Meanings and
Functions}, 1--32.
Amsterdam: John Benjamins.
\href{https://doi.org/10.1075/hcp.76.01sar}{doi:10.1075/hcp.76.01sar}.

\bibitem{kimoto2026existential}
Kimoto, Yukinori. 2026.
Variability in case-marking in northern Philippine languages:
Areal typology of existential and related clauses.
\textit{Gengo Kenkyu} 170: 1--27.
\href{https://doi.org/10.11435/gengo.170.0_1}{doi:10.11435/gengo.170.0\_1}.
